\chapter{Determinants}

The determinant of a square matrix is a single number
that summarises a great deal of geometric and algebraic information:
$\det A \neq 0$ iff $A$ is invertible iff the columns of $A$ are linearly
independent iff the system $A x = b$ has a unique solution for every $b$.
The determinant also gives the (signed) volume scaling factor of the linear
map $x \mapsto A x$. This chapter drills the computation in dimensions $2$
and $3$ and the structural properties.

\section{Core Rules}

\begin{rulebox}{The $2 \times 2$ formula}
For $A = \bigl(\begin{smallmatrix} a & b \\ c & d \end{smallmatrix}\bigr)$,
\[
  \det A \;=\; \begin{vmatrix} a & b \\ c & d \end{vmatrix} \;=\; a d - b c.
\]
The vertical-bar notation $\lvert A \rvert$ always denotes the determinant.
\end{rulebox}

\begin{rulebox}{The $3 \times 3$ formula (Sarrus' rule)}
For $A = \bigl(\begin{smallmatrix} a_{11} & a_{12} & a_{13} \\ a_{21} & a_{22} & a_{23} \\ a_{31} & a_{32} & a_{33} \end{smallmatrix}\bigr)$,
\[
  \det A \;=\; a_{11} a_{22} a_{33} + a_{12} a_{23} a_{31} + a_{13} a_{21} a_{32}
  - a_{13} a_{22} a_{31} - a_{11} a_{23} a_{32} - a_{12} a_{21} a_{33}.
\]
\end{rulebox}

\begin{rulebox}{Cofactor (Laplace) expansion}
Along row $i$:
\[
  \det A \;=\; \sum_{j = 1}^{n} (-1)^{i + j} \, a_{ij} \, M_{ij},
\]
where $M_{ij}$ is the determinant of the $(n - 1) \times (n - 1)$ submatrix
obtained by deleting row $i$ and column $j$. The signed minor
$(-1)^{i + j} M_{ij}$ is the cofactor $C_{ij}$.
\end{rulebox}

\begin{rulebox}{Key properties}
\begin{itemize}[leftmargin=1.6em,itemsep=0.3ex]
  \item Triangular shortcut: if $A$ is upper- or lower-triangular,
        $\det A = a_{11} a_{22} \cdots a_{nn}$ (product of the diagonal).
  \item $\det(I_{n}) = 1$.
  \item $\det(A^{T}) = \det A$.
  \item Multiplicativity: $\det(A B) = \det A \cdot \det B$.
  \item Scalar of a row: scaling \emph{one} row by $\lambda$ multiplies
        $\det A$ by $\lambda$. Hence $\det(\lambda A) = \lambda^{n} \det A$
        for $A \in \R^{n \times n}$.
  \item Row swap: swapping two rows flips the sign of $\det A$.
  \item Adding a multiple of one row to another leaves $\det A$ unchanged.
  \item A repeated row (or column) forces $\det A = 0$.
\end{itemize}
\end{rulebox}

\section{Worked Example}

\begin{worked}{Compute $\det M$ for $M=\begin{pmatrix}2&-1&3\\0&4&1\\5&2&-2\end{pmatrix}$ three ways, and settle what $\det$ does to $2M$ and to $M+N$}
\wstep{Step 1. Sarrus' rule, written out in full.}
For a $3\times3$ matrix only, copy the first two columns to the right and
take the three ``down-right'' products minus the three ``down-left'' ones:
\[
  \det M=\bigl[2\cdot4\cdot(-2)+(-1)\cdot1\cdot5+3\cdot0\cdot2\bigr]
        -\bigl[5\cdot4\cdot3+2\cdot1\cdot2+(-2)\cdot0\cdot(-1)\bigr]
\]
\[
  =\bigl[-16-5+0\bigr]-\bigl[60+4+0\bigr]=-21-64=\boxed{-85} .
\]
\wtrap{Sarrus works for $3\times3$ and \textbf{for nothing else}. There is no
$4\times4$ version, and inventing one is a guaranteed zero. For larger
matrices use cofactor expansion or row reduction.}

\wstep{Step 2. Cofactor expansion --- and choose the line with the zero in it.}
Expanding along the \emph{first column}, whose middle entry is $0$, kills
one of the three terms before it is written:
\[
  \det M=+2\begin{vmatrix}4&1\\2&-2\end{vmatrix}
        -0\begin{vmatrix}-1&3\\2&-2\end{vmatrix}
        +5\begin{vmatrix}-1&3\\4&1\end{vmatrix} .
\]
\[
  =2\bigl(4(-2)-1\cdot2\bigr)+5\bigl((-1)(1)-3\cdot4\bigr)
  =2(-10)+5(-13)=-20-65=-85 . \checkmark
\]
\wtrap{The signs alternate as
$\left(\begin{smallmatrix}+&-&+\\-&+&-\\+&-&+\end{smallmatrix}\right)$, starting from
$+$ at the top left --- the sign is $(-1)^{i+j}$ and belongs to the
\emph{position}, not to the entry. And each minor deletes the entry's whole
row \emph{and} whole column.}

\wstep{Step 3. Row reduction, which is what a computer does.}
$R_{3}\to R_{3}-\tfrac52R_{1}$ leaves the determinant unchanged:
\[
  \begin{pmatrix}2&-1&3\\0&4&1\\5&2&-2\end{pmatrix}
  \longrightarrow
  \begin{pmatrix}2&-1&3\\0&4&1\\0&\tfrac92&-\tfrac{19}{2}\end{pmatrix}
  \xrightarrow{\;R_{3}\to R_{3}-\tfrac98R_{2}\;}
  \begin{pmatrix}2&-1&3\\0&4&1\\0&0&-\tfrac{85}{8}\end{pmatrix}
\]
The determinant of a triangular matrix is the product of its diagonal:
$2\cdot4\cdot(-\tfrac{85}{8})=-85$. \checkmark\ Three methods, one answer.

\wstep{Step 4. What row operations do to a determinant.}
\begin{center}\footnotesize
\begin{tabular}{ll}
  \hline
  operation & effect on $\det$\\
  \hline
  add a multiple of one row to another & unchanged\\
  swap two rows & sign flips\\
  multiply a row by $c$ & multiplied by $c$\\
  \hline
\end{tabular}
\end{center}
Only the first is free; the other two must be recorded and undone at the
end.

\wstep{Step 5. Scaling the whole matrix: the factor is not what it looks like.}
Multiplying $M$ by $2$ multiplies \emph{every one of the three rows} by $2$,
so by the table above
\[
  \det(2M)=2^{3}\det M=8\cdot(-85)=-680 ,
\]
and in general $\det(cA)=c^{n}\det A$ for an $n\times n$ matrix.
\wtrap{$\det(2M)\neq2\det M$. The exponent is the \emph{size} of the matrix,
and this appears constantly in eigenvalue and definiteness work where a
matrix is rescaled.}

\wstep{Step 6. And addition: there is no rule at all.}
\[
  \det(A+B)\neq\det A+\det B .
\]
One counterexample settles it: $A=B=I_{2}$ gives $\det(A+B)=\det(2I)=4$
while $\det A+\det B=1+1=2$. By contrast the product \emph{does} behave:
\[
  \det(AB)=\det A\cdot\det B,
  \qquad
  \det(A^{-1})=\frac{1}{\det A},
  \qquad
  \det(A^{\!\top})=\det A .
\]

\wstep{Step 7. What the number $-85$ is for.}
It is non-zero, so $M$ is invertible, its rows are linearly independent,
$M\mathbf x=\mathbf b$ has exactly one solution for every $\mathbf b$, and
no eigenvalue is $0$. Every one of those statements is equivalent to
$\det M\neq0$; the determinant is a single number that decides all of them
at once. Its \emph{sign} and \emph{size} matter separately too: $|\det M|$
is the volume scaling factor of the map $\mathbf x\mapsto M\mathbf x$, and
the sign says whether orientation is reversed.
\end{worked}

\section{Practice Problems}

Evaluate each determinant, or answer the structural question posed.

\begin{problems}
  \item $\begin{vmatrix} 3 & 2 \\ 1 & 4 \end{vmatrix}$.
  \item $\begin{vmatrix} 5 & 7 \\ 2 & 3 \end{vmatrix}$.
  \item $\begin{vmatrix} 1 & 2 \\ 2 & 4 \end{vmatrix}$.
  \item $\begin{vmatrix} 6 & -4 \\ 3 & -2 \end{vmatrix}$.
  \item $\begin{vmatrix} \cos \theta & -\sin \theta \\ \sin \theta & \cos \theta \end{vmatrix}$.
  \item $\begin{vmatrix} 1 & 2 & 3 \\ 0 & 1 & 4 \\ 0 & 0 & 5 \end{vmatrix}$.
  \item $\begin{vmatrix} 2 & 0 & 0 \\ 5 & -1 & 0 \\ 7 & 3 & 4 \end{vmatrix}$.
  \item $\begin{vmatrix} 1 & 2 & 3 \\ 4 & 5 & 6 \\ 7 & 8 & 9 \end{vmatrix}$.
  \item $\begin{vmatrix} 1 & 0 & 2 \\ 3 & 1 & 0 \\ 0 & 4 & 5 \end{vmatrix}$.
  \item $\begin{vmatrix} 2 & -1 & 3 \\ 0 & 4 & 1 \\ 5 & 2 & 0 \end{vmatrix}$.
  \item $\begin{vmatrix} 1 & 1 & 1 \\ 1 & 2 & 3 \\ 1 & 4 & 9 \end{vmatrix}$ (Vandermonde-style).
  \item $\begin{vmatrix} 0 & 1 & 2 \\ 1 & 0 & 3 \\ 4 & 5 & 0 \end{vmatrix}$.
  \item Find all $x$ with $\begin{vmatrix} x & 2 \\ 3 & x + 1 \end{vmatrix} = 0$.
  \item Find all $\lambda$ with $\begin{vmatrix} 4 - \lambda & 1 \\ 2 & 3 - \lambda \end{vmatrix} = 0$.
  \item If $\det A = 5$ for $A \in \R^{3 \times 3}$, what is $\det(2 A)$?
  \item If $\det A = 5$ and $\det B = -2$ ($3 \times 3$), what is
        $\det(A B)$? $\det(A^{T} B)$?
  \item If $\det A = 3$ ($3 \times 3$), what is $\det(A^{-1})$? Justify.
  \item For an upper-triangular $4 \times 4$ matrix with diagonal entries
        $2, -1, 3, 5$, what is its determinant?
  \item Compute $\det A$ where
        $A = \begin{pmatrix} 2 & 4 & 6 \\ 1 & 3 & 5 \\ 0 & 0 & 7 \end{pmatrix}$.
  \item Compute $\begin{vmatrix} 1 & 2 & 0 \\ 3 & 6 & 1 \\ -1 & -2 & 4 \end{vmatrix}$.
  \item Use row operations (no Sarrus) to compute
        $\begin{vmatrix} 1 & 2 & 3 \\ 2 & 5 & 8 \\ 1 & 1 & 2 \end{vmatrix}$.
  \item If two rows of a $3 \times 3$ matrix $A$ are swapped to form $A'$,
        and $\det A = 7$, what is $\det A'$?
  \item For $A \in \R^{n \times n}$ with $\det A = 4$, give $\det(A^{5})$.
  \item Show that $\begin{vmatrix} a & b \\ b & a \end{vmatrix} = (a - b)(a + b)$.
  \item Compute $\begin{vmatrix} 1 & 0 & 0 & 0 \\ 2 & 3 & 0 & 0 \\ 4 & 5 & 6 & 0 \\ 7 & 8 & 9 & 10 \end{vmatrix}$.
\end{problems}

\section{Solutions}

\begin{solutions}
  \item $3 \cdot 4 - 2 \cdot 1 = 12 - 2 = 10$.
  \item $5 \cdot 3 - 7 \cdot 2 = 15 - 14 = 1$.
  \item $1 \cdot 4 - 2 \cdot 2 = 0$. (Rows proportional.)
  \item $6 \cdot (-2) - (-4) \cdot 3 = -12 + 12 = 0$.
  \item $\cos^{2} \theta - (-\sin \theta)(\sin \theta)
        = \cos^{2}\theta + \sin^{2}\theta = 1$.
  \item Upper-triangular: $1 \cdot 1 \cdot 5 = 5$.
  \item Lower-triangular: $2 \cdot (-1) \cdot 4 = -8$.
  \item Sarrus: $1 \cdot 5 \cdot 9 + 2 \cdot 6 \cdot 7 + 3 \cdot 4 \cdot 8
        - 3 \cdot 5 \cdot 7 - 1 \cdot 6 \cdot 8 - 2 \cdot 4 \cdot 9
        = 45 + 84 + 96 - 105 - 48 - 72 = 0$.
        (Row $3$ = $2 \cdot $ row $2 - $ row $1$.)
  \item Expand along row 1:
        $1 \cdot \begin{vmatrix} 1 & 0 \\ 4 & 5 \end{vmatrix}
        - 0
        + 2 \cdot \begin{vmatrix} 3 & 1 \\ 0 & 4 \end{vmatrix}
        = 1 \cdot 5 + 2 \cdot 12 = 5 + 24 = 29$.
  \item Expand along row 1:
        $2 \cdot \begin{vmatrix} 4 & 1 \\ 2 & 0 \end{vmatrix}
        - (-1) \cdot \begin{vmatrix} 0 & 1 \\ 5 & 0 \end{vmatrix}
        + 3 \cdot \begin{vmatrix} 0 & 4 \\ 5 & 2 \end{vmatrix}
        = 2 (0 - 2) + 1 (0 - 5) + 3(0 - 20) = -4 - 5 - 60 = -69$.
  \item Expand along row 1:
        $1 \cdot (2 \cdot 9 - 3 \cdot 4) - 1 \cdot (1 \cdot 9 - 3 \cdot 1)
        + 1 \cdot (1 \cdot 4 - 2 \cdot 1)
        = (18 - 12) - (9 - 3) + (4 - 2) = 6 - 6 + 2 = 2$.
  \item Expand along row 1:
        $0 - 1 \cdot (1 \cdot 0 - 3 \cdot 4) + 2 \cdot (1 \cdot 5 - 0 \cdot 4)
        = -1 \cdot (-12) + 2 \cdot 5 = 12 + 10 = 22$.
  \item $x(x + 1) - 6 = 0 \Rightarrow x^{2} + x - 6 = 0 \Rightarrow x = 2$
        or $x = -3$.
  \item $(4 - \lambda)(3 - \lambda) - 2 = 0 \Rightarrow
        \lambda^{2} - 7 \lambda + 10 = 0 \Rightarrow \lambda = 2$ or
        $\lambda = 5$.
  \item $\det(2 A) = 2^{3} \det A = 8 \cdot 5 = 40$.
  \item $\det(A B) = \det A \cdot \det B = -10$.
        $\det(A^{T} B) = \det(A^{T}) \det B = \det A \det B = -10$.
  \item $\det(A^{-1}) = 1/\det A = 1/3$. Reason: $A A^{-1} = I$ gives
        $\det A \cdot \det A^{-1} = 1$.
  \item $2 \cdot (-1) \cdot 3 \cdot 5 = -30$.
  \item Last row is $(0, 0, 7)$, so expand along it:
        $\det A = 7 \cdot \begin{vmatrix} 2 & 4 \\ 1 & 3 \end{vmatrix}
        = 7 \cdot 2 = 14$.
  \item Expand along row 1:
        $1 \cdot \begin{vmatrix} 6 & 1 \\ -2 & 4 \end{vmatrix}
        - 2 \cdot \begin{vmatrix} 3 & 1 \\ -1 & 4 \end{vmatrix}
        + 0 = (24 + 2) - 2(12 + 1) = 26 - 26 = 0$.
        (Col $2$ = $2 \cdot$ Col $1$.)
  \item Subtract $2 \cdot R_{1}$ from $R_{2}$ and $R_{1}$ from $R_{3}$:
        \[
          \begin{pmatrix} 1 & 2 & 3 \\ 0 & 1 & 2 \\ 0 & -1 & -1 \end{pmatrix}.
        \]
        Add $R_{2}$ to $R_{3}$:
        $\bigl(\begin{smallmatrix} 1 & 2 & 3 \\ 0 & 1 & 2 \\ 0 & 0 & 1 \end{smallmatrix}\bigr)$.
        Determinant $= 1 \cdot 1 \cdot 1 = 1$.
  \item Row swap flips sign: $\det A' = -7$.
  \item $\det(A^{5}) = (\det A)^{5} = 4^{5} = 1024$.
  \item $a \cdot a - b \cdot b = a^{2} - b^{2} = (a - b)(a + b)$.
  \item Lower-triangular: $1 \cdot 3 \cdot 6 \cdot 10 = 180$.
\end{solutions}

\section{Additional Problems}

\begin{problems}
  \item $\begin{vmatrix} 4 & 5 \\ 2 & 3 \end{vmatrix}$.
  \item $\begin{vmatrix} -2 & 3 \\ 4 & -1 \end{vmatrix}$.
  \item $\begin{vmatrix} 7 & 14 \\ 2 & 4 \end{vmatrix}$.
  \item $\begin{vmatrix} 6 & 1 \\ -3 & 2 \end{vmatrix}$.
  \item $\begin{vmatrix} 3 & 0 & 0 \\ 1 & 2 & 0 \\ 4 & 5 & 6 \end{vmatrix}$.
  \item $\begin{vmatrix} 2 & 1 & 3 \\ 0 & 4 & -1 \\ 0 & 0 & 5 \end{vmatrix}$.
  \item $\begin{vmatrix} 1 & 2 & 3 \\ 0 & 4 & 5 \\ 1 & 2 & 3 \end{vmatrix}$.
  \item $\begin{vmatrix} 1 & 2 & 1 \\ 3 & 0 & 2 \\ -1 & 4 & 1 \end{vmatrix}$.
  \item $\begin{vmatrix} 0 & 2 & 1 \\ 3 & 1 & 0 \\ 1 & 0 & 4 \end{vmatrix}$.
  \item $\begin{vmatrix} 2 & 1 & 1 \\ 1 & 2 & 1 \\ 1 & 1 & 2 \end{vmatrix}$.
  \item Find all $x$ with $\begin{vmatrix} x - 2 & 1 \\ 4 & x + 1 \end{vmatrix} = 0$.
  \item Find all $\lambda$ with $\begin{vmatrix} 5 - \lambda & 3 \\ 2 & 4 - \lambda \end{vmatrix} = 0$.
  \item If $\det A = -4$ for $A \in \R^{4 \times 4}$, what is $\det(3 A)$?
  \item If $\det A = 6$ and $\det B = 1/2$ ($3 \times 3$), what is
        $\det(A^{2} B^{T})$?
  \item If $\det A = 7$ ($n \times n$), what is $\det(A^{T} A)$?
  \item Compute $\begin{vmatrix} 2 & 1 & 3 \\ 1 & 0 & 2 \\ 4 & 2 & 6 \end{vmatrix}$
        (note relationship between rows).
  \item Use row operations to compute
        $\begin{vmatrix} 1 & 1 & 1 \\ 2 & 3 & 5 \\ 4 & 9 & 25 \end{vmatrix}$.
  \item For $A \in \R^{3 \times 3}$ with $\det A = 5$, what is $\det(-A)$?
  \item Show $\begin{vmatrix} a & b \\ -b & a \end{vmatrix} = a^{2} + b^{2}$.
  \item $\begin{vmatrix} 2 & 0 & 0 & 0 \\ 0 & -3 & 0 & 0 \\ 0 & 0 & 4 & 0 \\ 0 & 0 & 0 & -1 \end{vmatrix}$.
  \item $\begin{vmatrix} 1 & 0 & 0 & 5 \\ 0 & 2 & 0 & 0 \\ 0 & 0 & 3 & 0 \\ 0 & 0 & 0 & 4 \end{vmatrix}$.
  \item Find all $a$ such that
        $\begin{vmatrix} 1 & a & 0 \\ a & 1 & a \\ 0 & a & 1 \end{vmatrix} = 0$.
  \item Compute $\begin{vmatrix} 0 & 1 & 1 \\ 1 & 0 & 1 \\ 1 & 1 & 0 \end{vmatrix}$.
  \item For $A \in \R^{n \times n}$ with $\det A = 2$, what is
        $\det(A^{-1} A^{T})$?
  \item Use cofactor expansion along column 2 to compute
        $\begin{vmatrix} 3 & 0 & 2 \\ 5 & 1 & 4 \\ -1 & 0 & 3 \end{vmatrix}$.
\end{problems}

\section{Additional Solutions}

\begin{solutions}
  \item $4 \cdot 3 - 5 \cdot 2 = 12 - 10 = 2$.
  \item $(-2)(-1) - 3 \cdot 4 = 2 - 12 = -10$.
  \item $7 \cdot 4 - 14 \cdot 2 = 28 - 28 = 0$. (Rows proportional.)
  \item $6 \cdot 2 - 1 \cdot (-3) = 12 + 3 = 15$.
  \item Lower-triangular: $3 \cdot 2 \cdot 6 = 36$.
  \item Upper-triangular: $2 \cdot 4 \cdot 5 = 40$.
  \item Rows $1$ and $3$ equal: $\det = 0$.
  \item Expand along row 1:
        $1 \cdot (0 - 8) - 2 \cdot (3 + 2) + 1 \cdot (12 - 0)
        = -8 - 10 + 12 = -6$. Cross-check by Sarrus:
        $1 \cdot 0 \cdot 1 + 2 \cdot 2 \cdot (-1) + 1 \cdot 3 \cdot 4
        - 1 \cdot 0 \cdot (-1) - 1 \cdot 2 \cdot 4 - 2 \cdot 3 \cdot 1
        = 0 - 4 + 12 - 0 - 8 - 6 = -6$. Match.
  \item Expand along col 1:
        $-3 \cdot \begin{vmatrix} 2 & 1 \\ 0 & 4 \end{vmatrix}
        + 1 \cdot \begin{vmatrix} 2 & 1 \\ 1 & 0 \end{vmatrix}
        = -3 \cdot 8 + 1 \cdot (-1) = -24 - 1 = -25$.
        Cross-check by Sarrus:
        $0 \cdot 1 \cdot 4 + 2 \cdot 0 \cdot 1 + 1 \cdot 3 \cdot 0
        - 1 \cdot 1 \cdot 1 - 0 \cdot 0 \cdot 0 - 2 \cdot 3 \cdot 4
        = 0 + 0 + 0 - 1 - 0 - 24 = -25$. Match.
  \item Sarrus:
        $2 \cdot 2 \cdot 2 + 1 \cdot 1 \cdot 1 + 1 \cdot 1 \cdot 1
        - 1 \cdot 2 \cdot 1 - 2 \cdot 1 \cdot 1 - 1 \cdot 1 \cdot 2
        = 8 + 1 + 1 - 2 - 2 - 2 = 4$.
        Cross-check via row reduction:
        $R_{2} \to R_{2} - \tfrac{1}{2}R_{1}$,
        $R_{3} \to R_{3} - \tfrac{1}{2}R_{1}$ gives
        $\bigl(\begin{smallmatrix} 2 & 1 & 1 \\ 0 & 3/2 & 1/2 \\ 0 & 1/2 & 3/2 \end{smallmatrix}\bigr)$.
        Then $R_{3} \to R_{3} - \tfrac{1}{3}R_{2}$: $(0, 0, 4/3)$.
        $\det = 2 \cdot 3/2 \cdot 4/3 = 4$. Match.
  \item $(x - 2)(x + 1) - 4 = x^{2} - x - 2 - 4 = x^{2} - x - 6
        = (x - 3)(x + 2) = 0 \Rightarrow x = 3, -2$.
  \item $(5 - \lambda)(4 - \lambda) - 6 = \lambda^{2} - 9\lambda + 14
        = (\lambda - 2)(\lambda - 7) \Rightarrow \lambda = 2, 7$.
  \item $\det(3 A) = 3^{4} \det A = 81 \cdot (-4) = -324$.
  \item $\det(A^{2} B^{T}) = (\det A)^{2} \det B = 36 \cdot 1/2 = 18$.
  \item $\det(A^{T} A) = \det A^{T} \cdot \det A = (\det A)^{2} = 49$.
  \item Row 3 = $2 \cdot $ Row 1, so $\det = 0$. Verify by Sarrus:
        $2 \cdot 0 \cdot 6 + 1 \cdot 2 \cdot 4 + 3 \cdot 1 \cdot 2
        - 3 \cdot 0 \cdot 4 - 2 \cdot 2 \cdot 2 - 1 \cdot 1 \cdot 6
        = 0 + 8 + 6 - 0 - 8 - 6 = 0$. Match.
  \item $R_{2} \to R_{2} - 2 R_{1}$: $(0, 1, 3)$.
        $R_{3} \to R_{3} - 4 R_{1}$: $(0, 5, 21)$.
        $R_{3} \to R_{3} - 5 R_{2}$: $(0, 0, 6)$.
        $\det = 1 \cdot 1 \cdot 6 = 6$. Cross-check: Vandermonde with nodes
        $1, 2, 4$ gives $(2 - 1)(4 - 1)(4 - 2) = 1 \cdot 3 \cdot 2 = 6$.
        Match.
  \item $\det(-A) = (-1)^{3} \det A = -5$.
  \item $a \cdot a - b \cdot (-b) = a^{2} + b^{2}$.
  \item Diagonal: $2 \cdot (-3) \cdot 4 \cdot (-1) = 24$.
  \item Upper-triangular: $1 \cdot 2 \cdot 3 \cdot 4 = 24$. (The $5$ in
        position $(1, 4)$ is off-diagonal and doesn't affect.)
  \item Expand along row 1:
        $1 \cdot (1 - a^{2}) - a \cdot (a - 0) + 0
        = 1 - a^{2} - a^{2} = 1 - 2 a^{2} = 0
        \Rightarrow a = \pm 1/\sqrt{2}$.
  \item Sarrus:
        $0 + 1 \cdot 1 \cdot 1 + 1 \cdot 1 \cdot 1 - 1 \cdot 0 \cdot 1
        - 0 - 1 \cdot 1 \cdot 0 = 0 + 1 + 1 - 0 - 0 - 0 = 2$.
        Cross-check by expansion along row 1:
        $0 - 1 \cdot (0 - 1) + 1 \cdot (1 - 0) = 0 + 1 + 1 = 2$. Match.
  \item $\det(A^{-1} A^{T}) = \det(A^{-1}) \cdot \det(A^{T})
        = (1/\det A) \cdot \det A = 1$.
  \item Only the $(2, 2)$ entry in column 2 is nonzero.
        $\det = 1 \cdot (-1)^{2 + 2}
        \begin{vmatrix} 3 & 2 \\ -1 & 3 \end{vmatrix}
        = 1 \cdot (9 + 2) = 11$. Cross-check by Sarrus:
        $3 \cdot 1 \cdot 3 + 0 + 2 \cdot 5 \cdot 0 - 2 \cdot 1 \cdot (-1)
        - 3 \cdot 4 \cdot 0 - 0 = 9 + 0 + 0 + 2 - 0 - 0 = 11$. Match.
\end{solutions}
